% TOP_PROBLEM: {"id": 3312, "title": "Large induced forest in a planar graph.", "category_id": 3, "category": {"id": 3, "name": "graph_theory", "display_name": "Graph Theory", "description": "Problems involving graphs, networks, and their properties.", "slug": "graph-theory", "order_index": 3, "created_at": "2026-07-31T15:26:25.671Z"}, "status": "open", "problem_number": "OPG-46634", "set_id": 12}
% FIRST_POSTED: 2026-10-01
% ENTRY_KIND: statement_only
% UnsolvedMath ID 3312; source catalogue code OPG-46634
% !TeX program = lualatex
% Definition and sources only; this file is not an LLM solution attempt.
\documentclass[11pt,a4paper]{article}
\usepackage[margin=25mm]{geometry}
\usepackage{fontspec}
\setmainfont{lmroman10-regular.otf}[BoldFont=lmroman10-bold.otf,
 ItalicFont=lmroman10-italic.otf,BoldItalicFont=lmroman10-bolditalic.otf]
\usepackage{amsmath,amssymb,mathtools}
\usepackage{unicode-math}
\setmathfont{latinmodern-math.otf}
\AtBeginDocument{\let\setminus\smallsetminus}
\usepackage{xurl,hyperref}
\hypersetup{colorlinks=true,urlcolor=blue}
\setcounter{secnumdepth}{0}
\setlength{\parindent}{0pt}
\setlength{\parskip}{5pt}
\setlength{\emergencystretch}{3em}
\begin{document}
\section*{Large induced forest in a planar graph.}\label{3312}
{\small UnsolvedMath ID 3312 · OPG-46634 · Graph Theory · OpenGarden}
\subsection{Definitions and mathematical statement}
Let $G=(V,E)$ be a finite simple planar graph on $n$ vertices. For
$U\subseteq V$, its induced subgraph $G[U]$ has vertex set $U$ and
all edges of $G$ whose ends lie in $U$. A forest is an undirected
graph with no cycle; it may have several components and isolated
vertices. Define
\[
                      a(G)=\max\{|U|:G[U]\text{ is a forest}\}.
\]
The Albertson--Berman conjecture in this record is that
\[
                                a(G)\ge\frac n2
\]
for every finite simple planar graph. In integer terms this requires
at least $\lceil n/2\rceil$ vertices.

Equivalently, one should be able to delete a set $S\subseteq V$ of
at most $\lfloor n/2\rfloor$ vertices meeting every cycle. Such a set
is a feedback vertex set, and $G-S$ must retain every edge between
the remaining vertices. Selecting a spanning forest by simply
removing edges is therefore insufficient: a forest subgraph on all
vertices need not be induced.

The conclusion asks for one large induced forest, without requiring
that the complementary vertices also induce a forest. No girth,
maximum-degree, or connectedness assumption accompanies planarity.
The factor one half cannot be increased uniformly: a disjoint union
of copies of $K_4$ is planar and contributes at most two vertices
per copy to an induced forest, since every three vertices of a
$K_4$ already form a cycle. All bounds concern the number of
vertices rather than edges.
\subsection{Short English statement}
Can at least half of the vertices of every planar graph be retained while the induced graph has no cycle?
\subsection{Sources}
\begin{itemize}
\item[S1] ULAM.AI and the UnsolvedMath Contributors, supplied problems.json, record 3312 (OPG-46634), OpenGarden collection. Catalogue identity and source transcription; CC BY 4.0.
\url{https://huggingface.co/datasets/ulamai/UnsolvedMath}.
\item[S2] Open Problem Garden, Large induced forest in a planar graph.. Original problem and discussion; the mathematical formulation below is expanded from the supplied source record.
\url{http://www.openproblemgarden.org/op/large_induced_forest_in_a_planar_graph}.
\end{itemize}
\end{document}
